% \iffalse meta-comment
%
% Copyright (C) 2023 by Logan Grosz <logan.grosz@gmail.com>
% -------------------------------------------------------
% 
% This file may be distributed and/or modified under the
% conditions of the LaTeX Project Public License, either version 1.3
% of this license or (at your option) any later version.
% The latest version of this license is in:
%
%    http://www.latex-project.org/lppl.txt
%
% and version 1.3 or later is part of all distributions of LaTeX 
% version 2005/12/01 or later.
%
% \fi
%
% \iffalse
%<package>\NeedsTeXFormat{LaTeX2e}[2005/12/01]
%<package>\ProvidesPackage{roadsigns}
%<package>    [2023/05/14 v1.0 Graphics macros for drawing various roadsigns.]
%<package>\RequirePackage{tikz}
%<package>\usetikzlibrary{spath3,intersections}
%
%<*driver>
\ProvidesFile{roadsigns.dtx}
%</driver>
%
%<*driver>
\documentclass{ltxdoc}
\usepackage{roadsigns}
\EnableCrossrefs         
\CodelineIndex
\RecordChanges
\begin{document}
  \DocInput{roadsigns.dtx}
\end{document}
%</driver>
% \fi
%
% \CheckSum{162}
%
% \CharacterTable
%  {Upper-case    \A\B\C\D\E\F\G\H\I\J\K\L\M\N\O\P\Q\R\S\T\U\V\W\X\Y\Z
%   Lower-case    \a\b\c\d\e\f\g\h\i\j\k\l\m\n\o\p\q\r\s\t\u\v\w\x\y\z
%   Digits        \0\1\2\3\4\5\6\7\8\9
%   Exclamation   \!     Double quote  \"     Hash (number) \#
%   Dollar        \$     Percent       \%     Ampersand     \&
%   Acute accent  \'     Left paren    \(     Right paren   \)
%   Asterisk      \*     Plus          \+     Comma         \,
%   Minus         \-     Point         \.     Solidus       \/
%   Colon         \:     Semicolon     \;     Less than     \<
%   Equals        \=     Greater than  \>     Question mark \?
%   Commercial at \@     Left bracket  \[     Backslash     \\
%   Right bracket \]     Circumflex    \^     Underscore    \_
%   Grave accent  \`     Left brace    \{     Vertical bar  \|
%   Right brace   \}     Tilde         \~}
%
%
% \changes{v1.0}{2023/05/14}{Initial version}
% \changes{v1.1}{2023/08/24}{Added parking symbol}
%
% \GetFileInfo{roadsigns.sty}
%
% \title{The \textsf{roadsigns} package\thanks{This document
%   corresponds to \textsf{roadsigns}~\fileversion, dated \filedate.}}
% \author{Logan Grosz \\ \texttt{logan.grosz@gmail.com}}
%
% \maketitle
%
% \section{Introduction}
%
% This package provides several graphics macros for drawing
% roadsign graphics. It doesn't do much more than that.
%
% \section{Usage}
%
% Use any of the described macros anywhere a
% \textsf{tikzpicutre} environment could normally be used.
%
% \DescribeMacro{\interstate\marg{text}}
% This macro draws the U.S. interstate shield with
% \meta{text} as the number. The image only supports 2
% digit wide interstate numbers. The ``I" should not be
% included in \meta{text}.
%
% \DescribeMacro{\highway\marg{text}}
% This macro draws the generic U.S. highway shield with
% \meta{text} as the number. This image only supports 2
% digit highway numbers.
%
% \DescribeMacro{\widehighway\marg{text}}
% This macro draws the generic wide U.S highway shield with
% \meta{text} as the number. This image only supports 3
% digit highway numbers.
%
% \DescribeMacro{\statehighway\marg{text}}
% This macro draws a generic state highway ``pill." This
% image supports both 2 and 3 digit highway numbers.
%
% \DescribeMacro{\parking}
% This macro draws a parking sign with a simple ``P".
%
% \StopEventually{\PrintChanges}
%
% \section{Implementation}
%
% \begin{macro}{\interstate}
% Four circles are draw and split by their intersections. The
% components from the resulting split are combined in such
% a way to create the U.S. interstate shield shape.
%    \begin{macrocode}
\newcommand{\interstate}[1]{%
  \begin{tikzpicture}[x=1em,y=1em]
    \useasboundingbox (-0.8,-0.9) rectangle (0.8,0.8);
%    \end{macrocode}
% The following paths are framing for the shield.
%    \begin{macrocode}
    \path[spath/save=circleBR] (0.2,0.1) circle (0.9);
    \path[spath/save=circleBL] (-0.2,0.1) circle (0.9);
    \path[spath/save=circleTR] (0.3,1.5) circle (0.9);
    \path[spath/save=circleTL] (-0.3,1.5) circle (0.9);
    \tikzset
    {
%    \end{macrocode}
% Only split the paths at the appropriate points
%    \begin{macrocode}
      spath/split at intersections={circleTR}{circleTL},
      spath/split at intersections={circleTR}{circleBL},
      spath/split at intersections={circleBR}{circleTL},
      spath/split at intersections={circleBR}{circleBL},
%    \end{macrocode}
% Save the components for later
%    \begin{macrocode}
      spath/get components of={circleTR}\cTR,
      spath/get components of={circleBR}\cBR,
      spath/get components of={circleTL}\cTL,
      spath/get components of={circleBL}\cBL,
    }
%    \end{macrocode}
% Combines the appropriate components to make a shield
% shape
%    \begin{macrocode}
    \draw[
      fill=white,
      spath/use=\getComponentOf\cTR{4},
      spath/use={\getComponentOf\cBL{5},reverse,weld},
      spath/use={\getComponentOf\cBR{4},reverse,weld},
      spath/use={\getComponentOf\cTL{3},weld},
      spath/adjust and close=current,
    ];
    \node at (0,0) {\sffamily#1};
  \end{tikzpicture}
  }
%    \end{macrocode}
% \end{macro}
%
% \begin{macro}{\highway}
% 8 circles and two bars connecting the two circles with
% the highest vertical position are creating and split at
% appropriate intersections. The components from the
% resutling splits are combined to create the generic U.S.
% highway shield. In order to split on tangent circles, a
% line which is normal to their tangent point is drawn and
% the circles are split with that line.
%    \begin{macrocode}
\newcommand{\highway}[1]{%
  \begin{tikzpicture}[x=1em,y=1em]
    \useasboundingbox (-0.8,-0.8) rectangle (0.8,0.9);
%    \end{macrocode}
% Draws circles and lines which frame the generic U.S.
% highway shield
%    \begin{macrocode}
    \path[spath/save=circle1R] (0.217,1.077) circle (0.390);
    \path[spath/save=circle2R] (0.976,0.263) circle (0.40);
    \path[spath/save=circle3R] (0.35,-0.15) circle (0.35);
    \path[spath/save=circle4R] (0.35,-1.125) circle (0.625);
    \path[spath/save=connectorR] (0.217,1.077) -- (0.976,0.263);
    \path[spath/save=splitter1R] (0.976,0.263) -- (0.35,-0.15);
    \path[spath/save=splitter2R] (0.35,-0.15) -- (0.35,-1.125);
    \path[spath/save=circle1L] (-0.217,1.077) circle (0.390);
    \path[spath/save=circle2L] (-0.976,0.263) circle (0.40);
    \path[spath/save=circle3L] (-0.35,-0.15) circle (0.35);
    \path[spath/save=circle4L] (-0.35,-1.125) circle (0.625);
    \path[spath/save=connectorL] (-0.217,1.077) -- (-0.976,0.263);
    \path[spath/save=splitter1L] (-0.976,0.263) -- (-0.35,-0.15);
    \path[spath/save=splitter2L] (-0.35,-0.15) -- (-0.35,-1.125);
    \tikzset
    {
%    \end{macrocode}
% Splits the paths at appropriate intersections.
%    \begin{macrocode}
      spath/split at intersections={circle1R}{circle1L},
      spath/split at intersections={circle1R}{connectorR},
      spath/split at intersections={circle1L}{connectorL},
      spath/split at intersections={circle2R}{connectorR},
      spath/split at intersections={circle2L}{connectorL},
      spath/split at intersections={circle2R}{splitter1R},
      spath/split at intersections={circle2L}{splitter1L},
      spath/split at intersections={circle3R}{splitter1R},
      spath/split at intersections={circle3L}{splitter1L},
      spath/split at intersections={circle3R}{splitter2R},
      spath/split at intersections={circle3L}{splitter2L},
      spath/split at intersections={circle4R}{splitter2R},
      spath/split at intersections={circle4L}{splitter2L},
      spath/split at intersections={circle4R}{circle4L},
%    \end{macrocode}
% Save the components for use later
%    \begin{macrocode}
      spath/get components of={circle1R}\cOneR,
      spath/get components of={circle1L}\cOneL,
      spath/get components of={connectorR}\cConnectorR,
      spath/get components of={connectorL}\cConnectorL,
      spath/get components of={circle2R}\cTwoR,
      spath/get components of={circle2L}\cTwoL,
      spath/get components of={circle3R}\cThreeR,
      spath/get components of={circle3L}\cThreeL,
      spath/get components of={circle4R}\cFourR,
      spath/get components of={circle4L}\cFourL,
    }
%    \end{macrocode}
% Combines the appropriate components to make a shield
% shape.
%    \begin{macrocode}
    \draw[
      fill=white,
      spath/use={\getComponentOf\cOneR{3}},
      spath/use={\getComponentOf\cConnectorR{2},weld},
      spath/use={\getComponentOf\cTwoR{2},weld},
      spath/use={\getComponentOf\cThreeR{3},reverse,weld},
      spath/use={\getComponentOf\cFourR{2},weld},
      spath/use={\getComponentOf\cFourL{4},weld},
      spath/use={\getComponentOf\cThreeL{2},reverse,weld},
      spath/use={\getComponentOf\cTwoL{3},weld},
      spath/use={\getComponentOf\cConnectorL{2},reverse,weld},
      spath/use={\getComponentOf\cOneL{3},weld},
      spath/adjust and close=current,
    ];
    \node at (0,0) {\sffamily#1};
  \end{tikzpicture}
  }
%    \end{macrocode}
% \end{macro}
%
% \begin{macro}{\widehighway}
% Similar to \textsf{\\highway}. The circles were brought
% further apart apart and another line which is tangent to
% the bottom four circles is used to elongate the bottom of
% the shield.
%    \begin{macrocode}
\newcommand{\widehighway}[1]{%
  \begin{tikzpicture}[x=1em,y=1em,scale=1.2]
    \useasboundingbox (-1,-0.8) rectangle (1,0.8);
%    \end{macrocode}
% These are the circles which make up the curves of the
% shield. One circle was used to looks about correct
% (circle3R). The rests' positions were calculated by
% drawing circles over a real image and converting vectors
% in one span to another.
%    \begin{macrocode}
    \path[spath/save=circle1R] (0.3,1.11511) circle (0.522311);
    \path[spath/save=circle2R] (1.20532,0.190242) circle (0.38516);
    \path[spath/save=circle3R] (0.55,-0.15) circle (0.35);
    \path[spath/save=circle4R] (0.225,-0.794558) circle (0.294558);
    \path[spath/save=circle1L] (-0.3,1.11511) circle (0.522311);
    \path[spath/save=circle2L] (-1.20532,0.190242) circle (0.38516);
    \path[spath/save=circle3L] (-0.55,-0.15) circle (0.35);
    \path[spath/save=circle4L] (-0.225,-0.794558) circle (0.294558);
%    \end{macrocode}
% These lines connect the top two circles with an arbitrary line which looks ``correct."
%    \begin{macrocode}
    \path[spath/save=conn12R] (0.440624,0.867885) -- (1.1101,0.210911);
    \path[spath/save=conn12L] (-0.440624,0.867885) -- (-1.1101,0.210911);
%    \end{macrocode}
% This path is a long line tangent to the bottom four circles.
%    \begin{macrocode}
    \path[spath/save=conn34] (0.7,-0.50) -- (-0.7,-0.50);
%    \end{macrocode}
% These paths are used to split other paths which are not
% obviously intersecting, i.e. points where paths are
% tangent
%    \begin{macrocode}
    \path[spath/save=split23R] (1.20532,0.190242) -- (0.55,-0.15);
    \path[spath/save=split23L] (-1.20532,0.190242) -- (-0.55,-0.15);
    \path[spath/save=split3Conn34R] (0.55,-0.40) -- (0.55,-0.60);
    \path[spath/save=split4Conn34R] (0.225,-0.40) -- (0.225,-0.60);
    \path[spath/save=split3Conn34L] (-0.55,-0.40) -- (-0.55,-0.60);
    \path[spath/save=split4Conn34L] (-0.225,-0.40) -- (-0.225,-0.60);
    \tikzset
    {
%    \end{macrocode}
% Splits paths at appropriate intersections.
%    \begin{macrocode}
      spath/split at intersections={circle1R}{circle1L},
      spath/split at intersections={circle1R}{conn12R},
      spath/split at intersections={conn12R}{circle2R},
      spath/split at intersections with={circle2R}{split23R},
      spath/split at intersections with={circle3R}{split23R},
      spath/split at intersections with={circle3R}{split3Conn34R},
      spath/split at intersections with={conn34}{split3Conn34R},
      spath/split at intersections with={conn34}{split4Conn34R},
      spath/split at intersections with={circle4R}{split4Conn34R},
      spath/split at intersections={circle4R}{circle4L},
      spath/split at intersections={circle1L}{conn12L},
      spath/split at intersections={conn12L}{circle2L},
      spath/split at intersections with={circle2L}{split23L},
      spath/split at intersections with={circle3L}{split23L},
      spath/split at intersections with={circle3L}{split3Conn34L},
      spath/split at intersections with={conn34}{split3Conn34L},
      spath/split at intersections with={conn34}{split4Conn34L},
      spath/split at intersections with={circle4L}{split4Conn34L},
%    \end{macrocode}
% Saves the components for use later
%    \begin{macrocode}
      spath/get components of={circle1R}\circleOneR,
      spath/get components of={circle2R}\circleTwoR,
      spath/get components of={circle3R}\circleThreeR,
      spath/get components of={circle4R}\circleFourR,
      spath/get components of={circle1L}\circleOneL,
      spath/get components of={circle2L}\circleTwoL,
      spath/get components of={circle3L}\circleThreeL,
      spath/get components of={circle4L}\circleFourL,
      spath/get components of={conn12R}\connOneTwoR,
      spath/get components of={conn12L}\connOneTwoL,
      spath/get components of={conn34}\connThreeFour,
      spath/get components of={split23R}\splitTwoThreeR,
      spath/get components of={split23L}\splitTwoThreeL,
      spath/get components of={split3Conn34R}\splitThreeConnThreeFourR,
      spath/get components of={split4Conn34R}\splitFourConnThreeFourR,
      spath/get components of={split3Conn34L}\splitThreeConnThreeFourL,
      spath/get components of={split4Conn34L}\splitFourConnThreeFourL,
    }
%    \end{macrocode}
% Combines appropriate components to draw the shield shape.
%    \begin{macrocode}
    \draw[
      fill=white,
      spath/use={\getComponentOf\circleOneR{3}},
      spath/use={\getComponentOf\connOneTwoR{2},weld},
      spath/use={\getComponentOf\circleTwoR{2},weld},
      spath/use={\getComponentOf\circleThreeR{3},reverse,weld},
      spath/use={\getComponentOf\connThreeFour{2},weld},
      spath/use={\getComponentOf\circleFourR{2},weld},
      spath/use={\getComponentOf\circleFourL{2},weld},
      spath/use={\getComponentOf\connThreeFour{4},weld},
      spath/use={\getComponentOf\circleThreeL{2},reverse,weld},
      spath/use={\getComponentOf\circleTwoL{3},weld},
      spath/use={\getComponentOf\connOneTwoL{2},reverse,weld},
      spath/use={\getComponentOf\circleOneL{3},weld},
      spath/adjust and close=current,
    ];
    \node at (0,0) {\sffamily#1};
  \end{tikzpicture}
  }
%    \end{macrocode}
% \end{macro}
%
% \begin{macro}{\statehighway}
% A simple pill shape welded together.
%    \begin{macrocode}
\newcommand{\statehighway}[1]{%
  \begin{tikzpicture}[x=1em,y=1em]
%    \end{macrocode}
% These paths make up a pill - a rectangular shape capped
% on both horizontal ends by semicircles.
%    \begin{macrocode}
    \path[spath/save=top] (-0.5,0.6) -- (0.5,0.6);
    \path[spath/save=right] (0.5,0.6) arc (90:-90:0.6);
    \path[spath/save=bottom] (0.5,-0.6) -- (-0.5,-0.6);
    \path[spath/save=left] (-0.5,-0.6) arc (270:90:0.6);
%    \end{macrocode}
% Paths are combined and welded appropriately.
%    \begin{macrocode}
    \draw[
      fill=white,
      spath/use={top},
      spath/use={right,weld},
      spath/use={bottom,weld},
      spath/use={left,weld},
      spath/adjust and close=current,
    ];
    \node {\sffamily#1};
  \end{tikzpicture}
  }
%    \end{macrocode}
% \end{macro}
%
% \begin{macro}{\parking}
% A simple parking, ``P," sign.
%    \begin{macrocode}
\newcommand{\parking}{%
  \begin{tikzpicture}[x=1em,y=1em]
    \node at (0, 0) [rectangle,fill=white,draw=black,rounded corners=0.25em] {\sffamily P};
  \end{tikzpicture}
  }
%    \end{macrocode}
% \end{macro}
%
% \Finale
\endinput
